%%
%% CMPB-D-26-02177 — revised manuscript (R1), assembled 25 Oct 2026 from
%% _archive/submitted_2026-04/paper.tex and revision_2026-10/revision_passages_2026-10-21.tex.
%% Tables T1a/T1b/T2 are generated by scripts/make_revision_tables.py (tables/*.tex);
%% figures by scripts/make_revision_figures.py (figs/*.pdf).  Numbers: tables/numbers.md.
%% Template: cas-dc (double column).  <<QUERY: ...>> marks points for the authors.
%%

\documentclass[a4paper,fleqn]{cas-dc}

\usepackage[numbers]{natbib}
\usepackage[utf8]{inputenc}
\usepackage[T1]{fontenc}
\usepackage{amsmath,amssymb}
\usepackage{booktabs}
\usepackage{array}
\newcolumntype{L}[1]{>{\raggedright\arraybackslash}p{#1}}
\usepackage{siunitx}
\sisetup{output-decimal-marker={.}, group-separator={,}}
\usepackage{algorithm}
\usepackage{algpseudocode}
\usepackage{rotating}
\usepackage{lineno}
\linenumbers
\renewcommand\linenumberfont{\normalfont\tiny\sffamily}
\setlength\linenumbersep{3pt}

\graphicspath{{figs/}}
\emergencystretch=1em
% float placement: let tall single-column figures go on text pages or on float pages near their reference
\renewcommand{\topfraction}{0.95}
\renewcommand{\dbltopfraction}{0.95}
\renewcommand{\textfraction}{0.05}
\renewcommand{\floatpagefraction}{0.6}
\renewcommand{\dblfloatpagefraction}{0.6}
\setcounter{topnumber}{3}
\setcounter{dbltopnumber}{2}

\begin{document}
\let\WriteBookmarks\relax
\def\floatpagepagefraction{1}
\def\textpagefraction{.001}

% ─────────────────────────────────────────────────────────────
% META
% ─────────────────────────────────────────────────────────────

\shorttitle{Geodesic decay analysis of pace-mapping score maps}

\shortauthors{S.\ Willert et al.}

\title[mode = title]{Geodesic Decay Analysis: Quantifying the Spatial
  Resolution of Intracardiac Pace-Mapping Score Maps on Triangulated
  Atrial Surfaces}

% ─────────────────────────────────────────────────────────────
% AUTHORS (unchanged)
% ─────────────────────────────────────────────────────────────

\author[1]{Sven Willert}[orcid=0009-0008-0260-9721]
\cormark[1]
\ead{sw@willert.net}
\credit{Conceptualization, Methodology, Software, Formal Analysis,
  Writing -- Original Draft, Visualization}

\author[1]{Vera Maslova}[orcid=0009-0006-0647-767X]
\credit{Investigation}

\author[1]{Adrian Zaman}
\credit{Investigation}

\author[1]{Derk Frank}[orcid=0000-0001-7561-075X]
\credit{Supervision, Resources, Writing -- Review \& Editing}

\author[1]{Evgeny Lian}[orcid=0000-0001-9696-9227]
\credit{Investigation, Data Curation, Writing -- Review \& Editing, Validation}

\affiliation[1]{organization={University Hospital Schleswig-Holstein,
    Campus Kiel, Department of Internal Medicine III --
    Cardiology and Internal Intensive Care,
    Kiel University},
  addressline={Arnold-Heller-Stra\ss e 3, Haus C},
  city={Kiel},
  postcode={24105},
  country={Germany}}

\cortext[1]{Corresponding author}

% ─────────────────────────────────────────────────────────────
% ABSTRACT (structured; reader pass 27 Oct 2026; CMPB limit 350 words, four headings)
% ─────────────────────────────────────────────────────────────

\begin{abstract}
\textbf{Background and Objective:}
Intracardiac pace mapping displays the similarity between paced beats
and a template beat as a score map on the atrial surface, which is
judged visually as focal or diffuse.  Geodesic decay analysis replaces
this judgement by the width~$\sigma$ (mm) of the score decay along the
surface.
%
\textbf{Methods:}
From the vertex of maximum score, geodesic distances are computed with
the Heat Method, scores are averaged in 2-mm radial bins, and a
three-parameter Gaussian is fitted; $R^2$ and a focality index~$F$
flag inadequate fits and secondary high-score regions.  The estimator was
validated against analytic and exact geodesics, on synthetic Gaussian,
non-Gaussian and structured fields generated on 56 clinical atrial
meshes and on anisotropic fields on a sphere, and applied to 55~pace maps of 16~patients
(46 with a marked pacing site).
%
\textbf{Results:}
Heat-Method distances deviated from exact geodesics by $-1.3\,\%$
(median) within 30\,mm of the origin.  At 10\,\% independent vertex
noise the mean width bias was within $\pm 0.25$\,mm for 2--30\,mm and
the single-estimate SD 0.2--1.8\,mm for 3--30\,mm; correlated
deviations increased the scatter.  Clinically, $\sigma$ had a median of
7.6\,mm (IQR 4.8--15.6).  The width was associated with the distance
between the score maximum and the marked pacing site: the maximum lay
within 10\,mm in 25 of 30 maps with $\sigma < 10$\,mm, and the median
distance was 50\,mm in the seven maps with $\sigma \ge 20$\,mm.  A
single-chamber reference configuration was associated with larger
widths (exploratory; mixed-model
ratio 1.8, 95\,\% CI 1.04--3.09).
%
\textbf{Conclusions:}
Geodesic decay analysis measures the spatial resolution of a
pace-mapping score map, the width of its decay around the maximum,
reproducibly in millimetres and in under a second per map, with
fit-quality and focality flags to guide interpretation.
\end{abstract}

% KEYWORDS
\begin{keywords}
geodesic distance \sep triangulated surface \sep spatial resolution \sep
pace mapping \sep Heat Method \sep cardiac electrophysiology
\end{keywords}

\maketitle

% =============================================================
%  1  INTRODUCTION
% =============================================================
\section{Introduction}\label{sec:intro}

Electroanatomic mapping systems reconstruct three-dimensional surface
models of cardiac chambers and overlay colour-coded scalar fields ---
such as activation time, bipolar voltage, or matching score --- onto the
triangulated mesh~\citep{willert2024}.  Intracardiac pace mapping
compares the unipolar signal recorded at a reference catheter during
pacing at multiple sites with the signal recorded during a clinical
trigger~\citep{hayashi2016, lian2025} and projects the similarity score
onto the surface reconstruction as a heatmap.  A focal, sharply
delimited high-score region is commonly read as accurate localisation
of the trigger and a diffuse one as poor localisation, by eye.  The
present work turns this reading into a measurement and tests it against
the pacing sites marked during acquisition.

No established metric quantifies how sharply such a score map falls
off around its maximum, in physical units on the mesh surface.  This is
the sense in which we use spatial resolution: a property of the score
map, distinct from the localisation error of the mapping system, which
we examine separately against the marked pacing sites.  Prior electrophysiological
studies have treated resolution as a binary property --- whether two
pacing sites can be distinguished at a given
separation~\citep{kadish1991, azegami2005} --- and used Euclidean
distance, which underestimates the separation along the curved, folded
walls of the atria.  Geodesic distance, the shortest path along the
surface, is the appropriate metric but requires dedicated algorithms on
the triangulated mesh~\citep{grandits2024, williams2021}.

Characteristic spatial scales of fields on surfaces are estimated in
neighbouring literatures: the full width at half maximum of heat-kernel
or geodesic smoothing kernels on cortical surface
reconstructions~\citep{chung2005, hagler2006}, and the range of a
variogram or the decay of spatial autocorrelation~\citep{moran1950,
cressie1993}.  These methods typically describe the scale of a field
over a whole surface, often under an assumption of stationarity.  A
pace-mapping score map instead has one dominant maximum, and the quantity
of interest is the width of the decay around it.  The method presented
here is developed for and validated on this application; its transfer
to other peaked fields on anatomical surfaces is left open.

This paper makes three contributions:
%
\begin{enumerate}
  \item a \textbf{decay-width estimator} for score maps on triangulated
    atrial surfaces: Heat-Method geodesic distances from the score
    maximum, radial binning and a Gaussian width convention, with flags
    for inadequate fits ($R^2$) and secondary high-score regions ($F$);
  \item its \textbf{validation} against analytic and exact geodesics and
    on synthetic fields generated on the clinical meshes, which yields an
    operating range and the behaviour of the width under non-Gaussian,
    anisotropic and structured fields;
  \item a \textbf{clinical demonstration} on 55 pace maps of
    16~patients, in which the location of the score maximum is checked
    against the pacing site marked during acquisition.
\end{enumerate}
%
A width in millimetres allows score maps to be compared between regions
and patients on a stated scale instead of by impression, and at well
under a second per map it could be computed during the procedure.  The
flags state when a single radial profile does not describe a map, so
that such a map is inspected rather than summarised by one number.

\begin{figure*}[pos=!t]
  \centering
  \includegraphics[width=\textwidth]{fig_workflow.pdf}
  \caption{Geodesic decay analysis of one clinical pace map (map~28 of Supplementary Table~S12; left
    atrium, posterior wall, double-chamber reference; 6{,}812 vertices,
    mean edge length 1.9\,mm).  \textbf{a}~Triangulated surface with
    the vertex matching scores (colour bar in score points).
    \textbf{b}~Origin: the vertex of maximum score (star; Eq.~\ref{eq:origin}),
    with the mesh edges drawn; the triangle marks the pacing site that
    was marked during acquisition, 2\,mm away.  \textbf{c}~Geodesic
    distance from the origin computed with the Heat Method
    (Section~\ref{sec:heat_method}), with iso-distance lines at 10,
    20 and 30\,mm.  \textbf{d}~Radial binning: vertex scores against
    geodesic distance (grey; the horizontal bands reflect the integer
    scores reported by the system) and the means of the 2-mm bins
    (Eq.~\ref{eq:binning}).  \textbf{e}~Three-parameter Gaussian fit
    (Eq.~\ref{eq:gauss}) with amplitude~$A$, baseline~$B$ and
    width~$\sigma$.  \textbf{f}~Outputs: $\sigma$, the fit
    quality~$R^2$ and the focality index~$F$, computed on the set of
    vertices at or above the half maximum of the fitted profile (red:
    the component containing the origin; Eq.~\ref{eq:focality});
    $F$ is interpreted only when $R^2 \ge 0.7$.}
  \label{fig:workflow}
\end{figure*}

% =============================================================
%  2  DATA REPRESENTATION
% =============================================================
\section{Data Representation}\label{sec:data}

The chamber surface reconstructed by the mapping system is a triangle
mesh $(\mathcal{V}, \mathcal{T})$ with $N$~vertices
$\mathcal{V} = \{v_1, \ldots, v_N\}$ at positions
$\mathbf{p}_i \in \mathbb{R}^3$ and $M$~triangles~$\mathcal{T}$, a
piecewise-linear approximation of the atrial surface
$\mathcal{M} \subset \mathbb{R}^3$\label{sec:mesh}; of the largest
connected components of the 55 analysed meshes, 14 are closed and 41
have a boundary (Section~\ref{sec:cohort}).  A scalar field assigns a
value to each vertex\label{sec:scalar}:
%
\begin{equation}
  S \colon \mathcal{V} \to [0, 100], \quad v_i \mapsto s_i\,.
  \label{eq:scalar}
\end{equation}
%
In pace mapping, $s_i$ is the matching score displayed at the $i$-th
vertex: the percentage similarity between the signal recorded at the
reference catheter during pacing and the template beat.  The estimator
is defined for any vertex field with a single
dominant maximum; it is characterised here for matching-score maps.

% =============================================================
%  3  GEODESIC DISTANCE COMPUTATION
% =============================================================
\section{Geodesic Distance Computation}\label{sec:geodesic}

\subsection{Geodesic distance on surfaces}\label{sec:geodesic_def}

On a surface $\mathcal{M} \subset \mathbb{R}^3$ with induced metric
tensor~$g$, the geodesic distance between two points $p, q \in
\mathcal{M}$ is the infimum of the arc length over all connecting curves
on the surface:
%
\begin{equation}
  d_{\mathrm{geo}}(p, q) = \inf_{\gamma}
    \int_0^1 \!\sqrt{g_{\mu\nu}\,\dot{\gamma}^\mu(t)\,
    \dot{\gamma}^\nu(t)}\;\mathrm{d}t\,,
  \label{eq:geodesic}
\end{equation}
%
where the infimum is taken over all piecewise smooth curves
$\gamma\colon [0,1] \to \mathcal{M}$ with $\gamma(0)=p$ and
$\gamma(1)=q$.  Geodesic distance captures the true surface separation
along folds and curvatures and is in general longer than the Euclidean
(through-tissue) distance $\|\mathbf{p}_p - \mathbf{p}_q\|$ in
$\mathbb{R}^3$.

\subsection{Discrete algorithms}\label{sec:discrete_geo}

\emph{Dijkstra's algorithm} on the mesh graph~\citep{dijkstra1959} is
efficient ($\mathcal{O}(N\log N)$) but restricts paths to mesh edges and
overestimates distances (zigzag artefact; on our meshes by
\SI{5.7}{\percent} within 30\,mm of the origin, median map).
\emph{Fast Marching Methods}~\citep{kimmel1998} avoid this artefact but
are first-order accurate and sensitive to obtuse triangles.
\emph{Exact algorithms} such as MMP~\citep{mitchell1987, surazhsky2005}
give exact distances on the piecewise-flat surface at
$\mathcal{O}(N^2 \log N)$ (median 0.16\,s per source on our meshes,
against 0.03\,s for the Heat Method including its factorisation) and
serve as the reference of the validation (Section~\ref{sec:validation}).
The \emph{Heat Method}~\citep{crane2017} is close to exact at a cost
dominated by one sparse factorisation that is reused for every further
source, and was adopted here.

\subsection{The Heat Method}\label{sec:heat_method}

The Heat Method uses the relation between short-time heat diffusion
and geodesic distance~\citep{varadhan1967}: the normalised heat flux
points along the gradient of the distance function.  Distances are
obtained in three algebraic steps:
%
\begin{align}
  \bigl(M_{\mathrm{mass}} - t\,L_C\bigr)\,u &= \delta_{v^*}\,,
    \label{eq:heat1} \\[3pt]
  X &= -\,\nabla u \;/\; |\nabla u|\,,
    \label{eq:heat2} \\[3pt]
  \Delta\phi &= \nabla \cdot X\,,
    \label{eq:heat3}
\end{align}
%
where $M_{\mathrm{mass}}$ is the lumped mass matrix (the surface-area
share of each vertex), $L_C$ the cotangent Laplacian (the discrete
analogue of the second spatial derivative on the mesh; negative
semidefinite in the sign convention used here), $\delta_{v^*}$ the unit
impulse at the source vertex~$v^*$, defined in Eq.~\eqref{eq:origin} as
the vertex of maximum score, and $t$ a short time step proportional to
the square of the mean edge length.  Eq.~\eqref{eq:heat1} solves the heat
equation from the source; Eq.~\eqref{eq:heat2} extracts the unit
gradient direction of the distance function; and Eq.~\eqref{eq:heat3}
recovers the scalar distance field~$\phi$ via a Poisson solve, yielding
the geodesic distances $\hat{d}_i = \phi(v_i)$, with the additive
constant fixed so that the distance at the source is zero.  The heat
equation serves here only as a computational device; it does not model
electrical propagation.

\paragraph{Computational complexity and parameters.}
The two linear systems (Eqs.~\ref{eq:heat1} and~\ref{eq:heat3}) share
the sparsity pattern of the cotangent Laplacian and are prefactored
once by sparse Cholesky decomposition.  For matrices whose graph is
planar, as for a closed triangulated surface of genus zero, nested
dissection bounds the factorisation cost by $\mathcal{O}(N^{1.5})$
operations~\citep{george1973, lipton1979}; each subsequent single-source
query costs two forward and backward substitutions and one gradient
evaluation, i.e.\ $\mathcal{O}(\mathrm{nnz}(L))$.  On the 59 clinical
meshes ($N = 1{,}549$--$43{,}386$) the measured time of solver
construction scaled as $N^{1.19}$ and a query as $N^{1.11}$; the
largest mesh takes 0.19\,s to construct and 7\,ms per query (single
core; Supplementary Section~S16, with the hardware, the ordering and a
sparsity proxy of the factor).  Two settings enter the method, both left
at the library defaults for every result: the diffusion time
$t = \bar{h}^2$ ($\bar{h}$ the mean edge length, as recommended by
\citet{crane2017}) and the intrinsic-Delaunay (``robust'') Laplacian of
\citet{sharp2020}, which improves robustness to obtuse triangles.
No other setting of the distance solver is adjustable.

\paragraph{Implementation.}
Geodesic distances were computed with the C++ library
\texttt{geometry-central}~\citep{sharp2019} via its Python binding
\texttt{potpourri3d}~\citep{sharp2021} (version 1.4.0).

\subsection{Origin identification}\label{sec:origin}

The source vertex (origin of the decay analysis) is the vertex with the
highest scalar value:
%
\begin{equation}
  v^* = \arg\max_{v_i \in \mathcal{V}}\; s_i\,.
  \label{eq:origin}
\end{equation}
%
In pace mapping, $v^*$ is the site of best morphological match.  Several
vertices can share the maximum, because the system reports the score at
each pacing site as an integer percentage and synthetic fields saturate
at the upper score limit.  The origin is then the vertex nearest to the
centroid of the largest connected region of tied vertices (plateau
centre); a choice by vertex order would depend on the arbitrary
numbering of the mesh and would typically place the origin at the edge
of the plateau (rule and its effect on the clinical maps in
Supplementary Section~S19).  Because the maximum is taken over noisy
vertex values, the detected origin is displaced from the true source;
this displacement is part of the validation
(Section~\ref{sec:validation}).

% =============================================================
%  4  SPATIAL DECAY MODEL
% =============================================================
\section{Spatial Decay Model}\label{sec:decay}

\subsection{Radial binning}\label{sec:binning}

Given the geodesic distances $\hat{d}_i$ from the origin~$v^*$ to all
vertices, each pair $(\hat{d}_i, s_i)$ is assigned to a radial bin of
width $\Delta d = \SI{2}{\milli\metre}$.  The mean score in the $k$-th
bin is
%
\begin{equation}
  \bar{S}_k = \frac{1}{|\mathcal{B}_k|}
    \sum_{v_i \in \mathcal{B}_k} s_i\,,
  \label{eq:binning}
\end{equation}
%
where $\mathcal{B}_k = \{v_i : k\,\Delta d \le \hat{d}_i < (k{+}1)\,\Delta d\}$;
bins containing fewer than three vertices are excluded.  Because the
sampling density of the mesh is approximately uniform within local
neighbourhoods, the vertex mean approximates an area-weighted mean.  The bins are
weighted equally, and the fit is evaluated at the bin centres
$d_k = (k + \tfrac12)\,\Delta d$.  The
bin width of \SI{2}{\milli\metre} matches the mean edge length of the
analysed meshes (\SIrange{0.95}{2.78}{\milli\metre}) and still resolves
narrow peaks ($\sigma \le \SI{5}{\milli\metre}$); its sensitivity is
reported in Section~\ref{sec:disc_limitations}.

\subsection{Gaussian fit}\label{sec:gauss}

We model the decay of the scalar field as a function of geodesic
distance~$d$ from the origin:
%
\begin{equation}
  S(d) = A \exp\!\Bigl(-\frac{d^2}{2\sigma^2}\Bigr) + B\,,
  \label{eq:gauss}
\end{equation}
%
where $A$ is the amplitude above baseline (score points; fitted without a
sign constraint and positive in all clinical maps), $B$ the
baseline at large distance and $\sigma > 0$ the width parameter (mm).
The Gaussian is used as a \emph{width convention}, in the way the
Gaussian-equivalent full width at half maximum is used to report the
resolution of an imaging system, not as a physical model of the decay:
it has a single width parameter, a finite baseline, and a known relation
to the half-maximum width ($\mathrm{FWHM} = 2.355\,\sigma$).  A fixed
model also keeps the width comparable between maps, which a per-map
model choice would not.  The
consequences of other decay shapes for the width are quantified in
Sections~\ref{sec:validation_shape} and~\ref{sec:clinical_results}.

\paragraph{Reading the width.}
At $d = \sigma$ the fitted score has decayed to 61\,\% of the amplitude
above baseline, at $d = 2\sigma$ to 14\,\%.  Under the fitted model a
site at geodesic distance~$d$ from the maximum scores
$A\,[1 - \exp(-d^2/2\sigma^2)]$ less than the maximum, so $\sigma$ sets
the distance at which this difference becomes comparable to the scatter
of the scores.  The parameter
describes how sharply the score map localises, analogous to the width of
a point spread function; how it relates to the distance of the maximum
from the paced site is examined in Section~\ref{sec:clinical_localisation}.

\paragraph{Fit procedure.}
The three parameters are estimated by least-squares regression on the
binned profile:
%
\begin{equation}
  \mathcal{L}(A, \sigma, B) =
    \sum_{k}\bigl(\bar{S}_k - S(d_k;\,A,\sigma,B)\bigr)^2\,,
  \label{eq:cost}
\end{equation}
%
minimised with \texttt{scipy.optimize.curve\_fit}~\citep{virtanen2020},
using initial values $A_0 = \max_k \bar{S}_k - \min_k \bar{S}_k$,
$\sigma_0 = \SI{10}{\milli\metre}$, $B_0 = \min_k \bar{S}_k$, and the bound
$\sigma \in [0.1,\,200]\,\si{\milli\metre}$, which avoids artificial
truncation.

\paragraph{Fit quality flag.}
The coefficient of determination~$R^2$ of the Gaussian fit to the binned
profile serves as a model-adequacy flag.  Profiles with $R^2 < 0.7$ are
retained and reported with their~$\sigma$, but flagged as
\emph{inadequately described by a single radial Gaussian}.  The
threshold is a convention; its error rates on fields of known type are
given in Section~\ref{sec:validation_flags} and its effect on the
clinical maps in Section~\ref{sec:clinical_results}.


\subsection{Focality index}\label{sec:focality}

A single radial profile cannot represent a score field with more than
one region of high score.  To flag such fields we define a focality
index on the half-maximum set of the fitted profile.  Let
$\mathcal{H} = \{v_i : s_i \ge \hat B + \hat A/2\}$ be the set of
vertices whose score reaches at least the half maximum of the fitted
Gaussian (Eq.~\eqref{eq:gauss}), and let $C_0$ be the connected
component of~$\mathcal{H}$ on the mesh edge graph that contains the
origin~$v^*$.  The focality index is
%
\begin{equation}
  F = \frac{|C_0|}{|\mathcal{H}|} \in (0,\,1]\,.
  \label{eq:focality}
\end{equation}
%
$F$ counts vertices, so a separate region is weighted by its size
above the threshold, not by its height.  Maps with $F < 0.70$ are
flagged: at least 30\,\% of the vertices that reach the half maximum lie
in regions not connected to the peak at the origin.  Because
$\mathcal{H}$ depends on the fitted $\hat A$ and $\hat B$, $F$ is
interpreted only when $R^2 \ge 0.7$, and it is reported as not
available if $\mathcal{H}$ has fewer than three vertices or
$\hat A \le 0$.  $F$ is a screening flag; what it detects and misses is
characterised in Section~\ref{sec:validation_flags}.

\subsection{Pipeline summary and outputs}\label{sec:pipeline_summary}

Algorithm~\ref{alg:pipeline} summarises the pipeline; its runtime is
dominated by the factorisation, and for the largest mesh
($N = 43{,}386$) it completes in well under \SI{1}{\second}.  For every
map it returns the width~$\sigma$ with its class---below the validated
range ($\sigma < 2$\,mm), in range ($2 \le \sigma < 30$\,mm), broad
($30 \le \sigma < 50$\,mm), diffuse ($50 \le \sigma < 100$\,mm) or not
identifiable ($\sigma \ge 100$\,mm; Section~\ref{sec:validation_range})---the
fit quality~$R^2$ with its flag, the focality index~$F$ with its flag
(only when $R^2 \ge 0.7$), and the difference in AIC between the
Gaussian and an exponential fit,
$\Delta\mathrm{AIC} = \mathrm{AIC}_{\mathrm{Gauss}} - \mathrm{AIC}_{\mathrm{exp}}$
(positive when the exponential fits better).  The study analysis additionally computes an empirical half-maximum width
read from the binned profile (Supplementary Table~S12 and
Section~S20).  A flagged map is kept: the flag asks the reader to
inspect the profile or the map before using the number.

\begin{algorithm}[t]
\caption{Geodesic Decay Analysis}\label{alg:pipeline}
\small
\begin{algorithmic}[1]
\Require Mesh $(\mathcal{V}, \mathcal{T})$, scalar field $S$,
  bin width $\Delta d$
\Ensure $\sigma$ (with class), $R^2$ (flag), $F$ (flag)
\State $v^* \gets \arg\max_{v_i} s_i$; on ties, centre of the largest tied region
  \Comment{Eq.~\ref{eq:origin}}
\State Prefactor the heat-flow and Poisson systems (sparse Cholesky)
  \Comment{nested-dissection bound; Sec.~\ref{sec:heat_method}}
\State $\hat{d}_i \gets \Call{HeatMethod}{v^*}\ \forall\, v_i$
  \Comment{Eqs.~\ref{eq:heat1}--\ref{eq:heat3}}
\For{$k = 0, 1, 2, \ldots$}
  \State $\mathcal{B}_k \gets \{v_i : k\Delta d
    \le \hat{d}_i < (k{+}1)\Delta d\}$
  \If{$|\mathcal{B}_k| \ge 3$}
    $\bar{S}_k \gets |\mathcal{B}_k|^{-1}
      \sum_{v_i \in \mathcal{B}_k} s_i$
  \EndIf
\EndFor
\State Fit $S(d) = A e^{-d^2/2\sigma^2} + B$
  to $\{(d_k, \bar{S}_k)\}$; compute $R^2$
  \Comment{Eq.~\ref{eq:gauss}}
\State $\mathcal{H} \gets \{v_i : s_i \ge \hat B + \hat A/2\}$;
  $F \gets |C_0(\mathcal{H})| / |\mathcal{H}|$
  \Comment{Eq.~\ref{eq:focality}}
\State Flag $R^2 < 0.7$ (fit inadequate); if $R^2 \ge 0.7$, flag $F < 0.70$ (secondary region)
\State Class: below range if $\sigma < 2$, in range if $2 \le \sigma < 30$,
  broad if $30 \le \sigma < 50$, diffuse if $50 \le \sigma < 100$, not identifiable otherwise
\State \Return $\sigma, R^2, F$ with class and flags
\end{algorithmic}
\end{algorithm}

% =============================================================
%  5  SYNTHETIC VALIDATION
% =============================================================
\section{Synthetic Validation}\label{sec:validation}

\subsection{Design}\label{sec:validation_design}
The validation separates the three sources of error of the
estimator---the geodesic distances, the detection of the origin, and the
binned Gaussian fit---by generating synthetic score fields on distances
that do not come from the Heat Method.  Two kinds of ground truth were
used.  (i)~\emph{Analytic geodesics}: great-circle distances on a
sphere ($R = 25$\,mm, icosahedral triangulations with 642--10{,}242
vertices, mean edge 3.8--0.9\,mm) and unrolled-plane distances on open
cylinders ($R = 20$\,mm, height 80\,mm, 1{,}920--6{,}400 vertices,
regular and jittered triangulations).  (ii)~\emph{Exact polyhedral
geodesics} (Mitchell--Mount--Papadimitriou algorithm in the
implementation of \citet{surazhsky2005}, package \texttt{pygeodesic})
on all 56 non-fragmented clinical meshes (the 55 analysed maps and one
map excluded for quality; 40 distinct atrial reconstructions, because
maps of different regions of one chamber share the reconstruction),
from the vertex $v_0$ nearest to the centre of mass of the vertices, a
fixed source, independent of the clinical scores, whose anatomical
location varies between reconstructions.

A synthetic field was generated as
%
\begin{equation}
  S_{\mathrm{synth}}(v_i) = A\exp\!\Bigl(-\frac{d_i^{\,2}}
    {2\sigma_0^2}\Bigr) + B + \varepsilon_i\,,
  \label{eq:synth}
\end{equation}
%
with the ground-truth distance~$d_i$ from the source vertex, amplitude
$A = 35$, Gaussian vertex noise
$\varepsilon_i \sim \mathcal{N}(0, (\eta A)^2)$, clamping to
$[50, 100]$ and 5\,\% of vertices masked.  In the primary protocol the
baseline is $B = 55$, so that the noise-free peak (90) lies below the
upper score limit, as the clinical maxima mostly do
(Section~\ref{sec:clinical_results}).  Every field was processed by the
unchanged pipeline: maximum-score origin with the tie rule, Heat-Method
distances \emph{recomputed from the detected origin}, 2-mm binning,
Gaussian fit.  The primary protocol covered eleven true widths
$\sigma_0$ (2, 3, 5, 8, 10, 15, 20, 25, 30, 40 and 50\,mm) at noise
levels $\eta$ of 5, 10 and 15\,\% of the amplitude with 20 seeds per
mesh (36{,}960 fits; Table~\ref{tab:validation}, Fig.~\ref{fig:validation}).
A saturating sensitivity protocol with $B = 65$ places the noise-free
peak at the score limit (100), where noisy fields saturate in a plateau
around the source; it covers noise-free fields, known origins and widths
of 1--80\,mm, and the plateau-centre tie rule was chosen with it
(Supplementary Sections~S3 and~S4).

Independent vertex noise is a reference condition: the within-bin
scatter of the clinical fields is larger (median 36\,\% of the fitted
amplitude, IQR 21--46\,\%).  Spatially structured deviations, which are
candidate sources of this excess, are examined in
Section~\ref{sec:validation_shape}.

\subsection{Results}\label{sec:validation_results}
\paragraph{Geodesic distances.}
On the analytic surfaces, where the geodesic is known in closed form,
the exact polyhedral solver was within 0.3\,\% of it, the Heat Method
underestimated it by 0.6--1.7\,\% and Dijkstra's algorithm
overestimated it by 6--18\,\% (Supplementary Table~S1, Fig.~S1).  On
the 56 clinical meshes, with the exact solver as reference, the
Heat-Method distance deviated by a median of $-1.3$\,\% (IQR $-1.6$ to
$-1.1$\,\%) between 2 and 30\,mm from the origin, where the width is
determined, and by $-0.2$\,\% over the whole mesh (Supplementary
Table~S2).  On Gaussian fields generated with
the analytic distances, Heat-Method distances lowered the recovered
width by a median of 0.1--0.2\,mm and Dijkstra distances widened it by
0.6--1.1\,mm (Supplementary Table~S1b).

\paragraph{Width recovery.}
In the primary protocol at 10\,\% noise, the mean bias was within
$\pm 0.1$\,mm for $\sigma_0 = 3$--30\,mm, and the SD of single estimates
increased from 0.23 to 1.75\,mm (Table~\ref{tab:validation},
Fig.~\ref{fig:validation}).  The detected origin lay a median of
0.6--6.9\,mm from the true source.  At 5 and 15\,\% noise, the bias
remained small over the same width range (Supplementary Table~S4b).
At $\sigma_0 = 2$\,mm and 10\,\% noise, 90\,\% of the estimates lay
between 1.7 and 2.5\,mm, but one runaway fit increased the SD to
5.9\,mm (Table~\ref{tab:validation}; Supplementary Section~S4b).
At 40 and 50\,mm and 10\,\% noise, the mean bias remained within
$\pm 0.4$\,mm, while the SD increased to 6.0 and 12.4\,mm.

In the saturating sensitivity protocol (Supplementary Tables~S3, S3b
and S4) the mean bias without noise was within $\pm 0.09$\,mm for
$\sigma_0 = 3$--30\,mm and, with the detected origin at 10\,\% noise,
between 0.00 and $+0.21$\,mm for 2--30\,mm (SD at most 1.24\,mm).
Below 2\,mm the estimator was unstable; at 60 and 80\,mm the SD reached
13 and 31\,mm.

\input{tables/T1a_validation_iid}


\begin{figure}[pos=!t]
  \centering
  \includegraphics[width=\linewidth,height=0.66\textheight,keepaspectratio]{fig_validation.pdf}
  \caption{Synthetic validation on the 56 clinical meshes with
    exact-geodesic ground truth, primary protocol (noise-free peak below
    the score limit, origin detected by the pipeline; $n = 1{,}120$ fits
    per point).  \textbf{A}~Bias $\hat\sigma - \sigma_0$ (mean, 95\,\%
    bootstrap band) for vertex noise of 5, 10 and 15\,\% of the
    amplitude.  \textbf{B}~Standard deviation of $\hat\sigma$.
    \textbf{C}~Operating range at 10\,\% noise over
    $\sigma_0 = 2$--50\,mm: median bias with the 5th--95th percentile
    band; shaded: the classes of Section~\ref{sec:validation_range}.
    Noise-free and known-origin results, the saturating sensitivity
    protocol and
    widths of 1--1.5 and 60--80\,mm are in Supplementary Tables~S3, S3b
    and S4.}
  \label{fig:validation}
\end{figure}

\subsection{Model deviations}\label{sec:validation_shape}
\paragraph{Non-Gaussian decay.}
For exponential and Lorentzian profiles of known half-maximum
width (7--47\,mm; 17{,}920 fits on the 56 meshes) the Gaussian-equivalent
width $2.355\,\hat\sigma$ overestimated it by 25--64\,\% and 6--34\,\% at
10\,\% noise; for Gaussian and raised-cosine profiles it was within $-8$
to $+1$\,\%.  An exponential fit had the lower AIC in at least 96\,\% of
the exponential and Lorentzian fields except the broadest Lorentzian
(Supplementary Section~S6), so the pipeline reports
$\Delta\mathrm{AIC}$ per map.
\paragraph{Anisotropic decay.}
For elliptical Gaussian fields with axis ratios up to four the isotropic
width approximated the geometric mean of the principal widths (ratio
0.87--1.04 at 10\,\% noise) with a median $R^2 \ge 0.97$: the isotropic
width is an azimuthal average, and $R^2$ remained high in these
anisotropic simulations
(Supplementary Section~S7).
\paragraph{Structured deviations.}
Spatially correlated noise of 0.25\,$A$ (smoothing scale 5 and 10\,mm)
displaced the detected origin by more than 10\,mm in 17--34\,\% of the
fields; among trials with the detected origin within 10\,mm of the true
source, the median width bias was within $\pm 0.7$\,mm with an SD of
1.8--3.5\,mm (Supplementary
Section~S9, Table~S9b).  Structured deviations thus cause two failure
modes: capture of the origin by a secondary maximum, which the flags
detect only in part, and a lowered $R^2$ with a more scattered width.
They also raise the within-bin scatter (15--19\,\% of~$A$ at 0.25\,$A$,
against 36\,\% in the clinical maps; Supplementary Table~S10).
\paragraph{Origin rule.}
An alternative origin, the vertex with the highest mean score in its
3-mm neighbourhood, did not reduce the synthetic bias and lay further
from the marked pacing site in the clinical maps (median 10.1 vs
6.3\,mm; Supplementary Section~S5).

\subsection{Classification flags}\label{sec:validation_flags}
On the sphere, fields of known type (Gaussian, flat, ring-shaped and
two-peaked; vertex noise 5--25\,\%, 20 trials per condition) were
classified with the two flags.  The $R^2$ flag rejected no Gaussian field
at thresholds up to 0.8 and accepted flat fields as adequate in at most
15\,\% of the trials; ring-shaped and two-peaked fields pass, because
their radial average still decays.  The $F$ flag detected two peaks at
least 25\,mm apart with an amplitude ratio of 0.7 or more in
90--100\,\% of the trials at 5--10\,\% noise; peaks of half the
amplitude and rings were not detected, and peaks 15\,mm or less apart
were usually missed (detection 0--15\,\%).  It flagged a
single peak in at most 2\,\% of the trials under independent noise up to
15\,\%, in 90\,\% at 25\,\% noise, and in 45--69\,\% of the fields with
correlated deviations of 0.25\,$A$ on the meshes (Supplementary
Section~S11, Tables~S11c and~S14b).

\subsection{Operating range}\label{sec:validation_range}
For Gaussian fields with independent vertex noise up to 10\,\% the mean
bias is within $\pm 0.25$\,mm for $\sigma_0 = 2$--30\,mm in both
protocols, and the SD of single estimates is 0.2--1.8\,mm for 3--30\,mm
(class ``in range'', $2 \le \sigma < 30$\,mm); 2\,mm is an unstable lower
edge with rare runaway fits.  Between 30 and 50\,mm the bias is small,
but the SD reaches 6--12\,mm in the primary protocol (class
``broad''), so such widths are
compared only coarsely.  Widths of 50\,mm and above are reported as
diffuse and not interpreted, fits at or above 100\,mm as not
identifiable.  The class of a map is assigned from its estimate
$\hat\sigma$, which need not fall in the class of the true width: a
broad field can yield a small estimate (Supplementary Table~S4).  Under
correlated deviations of 0.25\,$A$ the SD of the width rose to
2.5--9.1\,mm (1.8--3.5\,mm with the origin within 10\,mm of the true
source; Section~\ref{sec:validation_shape}), so the range above is a
best case.

For two independent estimates the SD of the difference is
$\sqrt{s_1^2 + s_2^2}$.  At 10\,\% independent noise a difference
exceeds the estimator's variability (approximately 95\,\%) when it is
larger than about 0.8\,mm at 5\,mm, 1.3\,mm at 10\,mm, 1.8\,mm at
15\,mm, 2.7\,mm at 20\,mm and 4.9\,mm at 30\,mm.  These are illustrative
scales from SDs pooled over meshes and noise realisations, not the
repeatability of an individual clinical map.


% =============================================================
%  6  CLINICAL DEMONSTRATION
% =============================================================
\section{Clinical Demonstration}\label{sec:clinical}

\subsection{Data}\label{sec:cohort}

Fifty-nine pace maps from 16~patients were acquired with the CARTO-3
electroanatomic mapping system (Biosense Webster) in four anatomic
regions of the left and right atrium (LA~posterior, LA~anterior,
RA~posterior, RA~lateral) at the University Hospital
Schleswig-Holstein, Kiel.  Three maps with fragmented meshes and one map
with 12 pacing sites and a maximum score of 68 were excluded, leaving
55~maps for the decay analysis (1{,}549--43{,}386 vertices, mean edge
length \SIrange{0.95}{2.78}{\milli\metre}).  In each region, pacing was
performed sequentially from several sites at a constant cycle length
(600\,ms); the first paced beat of a map served as the template, so
that the matching score of every later site is its similarity to a
paced beat (simulated trigger).  The analysis uses the score map as
exported by the mapping system, with one matching score per mesh vertex;
the maps comprised 7--135 pacing sites (median 48).  Eight patients were
mapped with a single-chamber reference, a decapolar catheter in the
coronary sinus, and eight with a double-chamber reference, following the
biatrial reference concept of \citet{lian2025}, adapted here: six
electrodes of the coronary-sinus catheter and four electrodes of a
right-atrial catheter are combined through an external pin box so that
the mapping system registers them as one ten-pole reference catheter.
No patient was mapped with both, so the reference configuration is a
between-patient factor, whereas the four regions are within-patient (one
to five maps per patient, median four).
% catheter description follows the thesis Methods (thesis-med/_mainmatter/methoden.tex); Lian et al. 2025 (EnSite X) used 3 CS + 3 HRA electrodes via an adapter to the precordial leads -> cited for the concept only
During acquisition, the operator marked the pacing site whose beat
served as the template with a unique marker in 46 of the 55 maps (a
``Location Only'' tag; the site scores 100 against its own beat).  The
nine maps without a unique marker are left out of the localisation
check (Section~\ref{sec:clinical_localisation}).

Mesh and score data were imported from the CARTO-3 system with
pulse-ep~\citep{willert2026pulseep}, our open-source platform for
electroanatomical mapping data, and the imported
geometry, scalar fields and colour coding were checked visually against
the CARTO display for each map.

The study was conducted in accordance with the Declaration of Helsinki
and approved by the institutional ethics committee of the Medical Faculty
of Christian-Albrechts-Universit\"at zu Kiel (protocol no.\ D\,418/23).
All patients provided written informed consent for the
electrophysiological procedure and for the use of anonymised mapping data
in research.

\subsection{Results}\label{sec:clinical_results}
\paragraph{Widths and fit quality.}
The Gaussian fit converged for all 55 maps (per-map results in
Supplementary Table~S12).  The width had a median of 7.6\,mm (IQR
4.8--15.6\,mm); 46 maps were in range, four broad, four diffuse and one
below the range (Table~\ref{tab:clinical_groups}).  The maximum vertex
score had a median of 95 and reached the upper limit of 100 in two maps.
The median~$R^2$ was 0.82, and 13 maps were flagged ($R^2 < 0.7$;
thresholds of 0.6 and 0.8 flag 7 and 23).  The flagged maps were mostly
narrow (median width 4.8\,mm; one broad map, $\sigma = 36.3$\,mm) with a smaller amplitude above baseline
(median 18 vs 26 score points), and a median of 98\,\% of their squared
residual lay beyond $3\sigma$: a low $R^2$ marked a narrow peak over an
irregular or secondarily elevated far field
(Fig.~\ref{fig:clinical_example}A).

\begin{figure}[pos=!tp]
  \centering
  \includegraphics[width=\linewidth,height=0.78\textheight,keepaspectratio]{fig_heatmap.pdf}
  \caption{Matching-score maps on the atrial surface (sequential colour
    map in score points; view along the surface normal at the origin;
    star: vertex of maximum score; triangle: pacing site marked during
    acquisition; lines: geodesic distance $\sigma$ and $2\sigma$ from the
    origin; scale bar 10\,mm; map numbers as in Supplementary Table~S12).
    \textbf{A:}~Map~75, LA posterior wall, single-chamber reference
    ($\sigma = 6.8$\,mm, $R^2 = 0.82$).
    \textbf{B:}~Map~78, RA lateral wall, single-chamber reference
    ($\sigma = 34.3$\,mm, $R^2 = 0.98$; broad): the $\sigma$ contour
    encloses most of the visible wall and the marked site lies 55\,mm
    from the maximum.}
  \label{fig:heatmap}
\end{figure}

\begin{figure*}[pos=!tp]
  \centering
  \includegraphics[width=\textwidth,height=0.66\textheight,keepaspectratio]{fig_clinical_example.pdf}
  \caption{Two flagged clinical maps with intermediate results (the
    unflagged map~28 is shown in Fig.~\ref{fig:workflow}).  Columns:
    score map with the origin (star), the marked pacing site (triangle)
    and the $\sigma$ and $2\sigma$ contours; geodesic distance from the
    origin; radial profile with 2-mm bin means and Gaussian fit (dashed
    when $R^2 < 0.7$; fitted over the full range, first 40\,mm shown; inset in A: full range); the half-maximum set of the fitted profile (red:
    component containing the origin; blue: other components) with the
    outputs; map numbers as in Supplementary Table~S12.
    \textbf{A:}~Narrow peak over a structured far field (map~1), LA anterior
    wall, double-chamber reference ($\sigma = 4.8$\,mm, $R^2 = 0.47$):
    the fit is flagged; $F = 0.45$ is not interpreted because the
    half-maximum set depends on the inadequate fit.
    \textbf{B:}~Second region above the half maximum (map~25), RA lateral wall,
    double-chamber reference ($\sigma = 4.6$\,mm, $R^2 = 0.97$,
    $F = 0.27$ with six components): the peak at the origin is narrow
    and well fitted, but regions 10--25\,mm away reach 78\,\% of its
    height; the map is flagged by~$F$.}
  \label{fig:clinical_example}
\end{figure*}

\input{tables/T2_clinical_groups}

\paragraph{Shape of the decay.}
On the binned profiles an exponential model had the lower AIC in 29
maps ($\Delta\mathrm{AIC} > 2$), the Gaussian in 15, and 11 were
undecided; the differences lay mainly in the far field (median RMSE
within $3\sigma$ 1.9 score points for both models; Supplementary
Section~S20).  Map rankings were strongly
concordant across the two models (Spearman $\rho = 0.97$ between the
Gaussian~$\sigma$ and the exponential scale); the widths in millimetres
depend on the model.
\paragraph{Distance definition.}
Recomputing the width from the same origin with exact geodesics changed
it by a median of $+0.08$\,mm, and Dijkstra distances by $+0.49$\,mm.
Within 30\,mm of the origin the Euclidean chord was shorter than the
geodesic by a median of 3.8\,\% (95th percentile within maps 11\,\%).
The Euclidean width differed from the geodesic width by a median of
$-0.29$\,mm, and the two were strongly rank-correlated (Spearman
$\rho = 0.94$), but the difference was concentrated in some maps: it
exceeded 3\,mm in 16, mostly broad or folded (8 of the 35 adequately fitted maps with
$\sigma < 30$\,mm; e.g.\ 11.0 vs 67.5\,mm in a lateral right-atrial
map).  The Euclidean fit flagged 17 maps compared with 13 for geodesic
distances; seven maps gained the flag and three lost it (Supplementary
Section~S13).
\paragraph{Focality index.}
Of the 42 maps with an adequate fit, four were flagged by $F$ (all
double-chamber): narrow peaks ($\sigma$ = 2.3--4.6\,mm) with further
regions above the half maximum, the nearest 6--18\,mm from the origin
(Fig.~\ref{fig:clinical_example}B).  $F$ had a median of 1.00 (IQR
0.89--1.00); thresholds of 0.5 to 0.9 flagged 3 to 11 of the 42 maps,
and $F$ was related to~$\sigma$ ($\rho = 0.51$; narrow peaks were more
often accompanied by comparably high regions elsewhere) and weakly
to~$R^2$ ($\rho = -0.28$; Supplementary Table~S14).

\subsection{Width and location of the maximum}\label{sec:clinical_localisation}
\begin{figure}[pos=!tp]
  \centering
  \includegraphics[width=\linewidth]{fig_marker_check.pdf}
  \caption{Geodesic distance from the vertex of maximum score to the
    pacing site marked during acquisition against the width~$\sigma$
    (46 maps; colour: fit-quality flag; logarithmic width axis, linear
    distance axis; the five narrow maps with a distance above 10\,mm are
    labelled with their map number; top: number of maps within 10\,mm or
    median distance for widths below 10, from 10 to 20 and from 20\,mm).
    Dashed: 10-mm distance; dotted:
    $\sigma = 10$\,mm, the model-based reading of
    Section~\ref{sec:disc_comparison}.}
  \label{fig:marker_check}
\end{figure}

In the 46 maps with a marked pacing site, the geodesic distance from
the vertex of maximum score to the marked site had a median of 6.3\,mm
(IQR 2.6--23.9\,mm) and increased with the width (Spearman $\rho = 0.51$
between offset and width, patient-cluster bootstrap 95\,\% interval
0.23--0.72; with $R^2$, $\rho = -0.05$; Fig.~\ref{fig:marker_check}).
The 30 maps with $\sigma < 10$\,mm came from 13 patients; nine carried
the fit-quality flag, four the focality flag, and one had a width below
2\,mm.  In these maps the distance was 4.1\,mm (IQR 2.2--7.2\,mm), and the
maximum lay within 10\,mm of the marked site in 25 of them; in the 9
maps with $10 \le \sigma < 20$\,mm it was 6.7\,mm, and in the 7 maps with
$\sigma \ge 20$\,mm 50\,mm (IQR 39--54\,mm).  In five of the 30 narrow
maps the maximum lay 20--27\,mm from the marked site, i.e.\ another
location of the map reached the highest score; neither flag identified
these maps reliably ($R^2$ flagged one, $F$ one).  Offsets are measured
to the mesh vertex nearest to the marked position (median distance
2.1\,mm, maximum 6.1\,mm).  Using the marked site as origin gave similar median widths when the
origins were within 5\,mm, whereas larger differences occurred when
they were more than 10\,mm apart (Supplementary Section~S18).  The
pipeline keeps the maximum as origin and uses the marked site as an
external reference.

\subsection{Reference configuration and region (exploratory)}\label{sec:clinical_reference}
Widths were larger with a single-chamber than with a double-chamber
reference (map-level medians 10.4 vs 5.2\,mm;
Table~\ref{tab:clinical_groups}, Supplementary Fig.~S6).  With
the patient as the unit of correlation, a linear mixed model of
$\log\sigma$ with reference and region as fixed effects and a patient
random intercept gave a single-/double-chamber width ratio (ratio of
geometric means) of 1.79 (95\,\% CI 1.04--3.09) and a ratio of 1.73
(0.85--3.53) for the lateral right-atrial wall relative to the posterior
left-atrial wall; the patient random intercept was small (intraclass
correlation 0.03).  At the patient level (8 vs 8) the difference was less
clear: a two-stage analysis gave a ratio of 1.75 (0.93--3.29;
diagnostics and further estimands in Supplementary Section~S15).  With
eight patients per configuration, neither configuration nor patients
randomised, these estimates describe an association of low precision
rather than an effect of the reference; the clinical analysis of the
configurations will be reported separately.


% =============================================================
%  7  DISCUSSION
% =============================================================
\section{Discussion}\label{sec:discussion}

\subsection{What the width adds}\label{sec:disc_method}
Geodesic decay analysis reduces a pace-mapping score map to one number,
the width~$\sigma$ of its decay along the atrial surface, with flags for
maps that a single radial profile does not describe.  It is simple (one
geodesic distance field, a binned profile and a three-parameter fit),
fast (well under a second per map) and read in millimetres.  The
validation shows that its distances reproduce exact geodesics to about
1\,\% where the width is determined (width difference $+0.08$\,mm), that
its bias is small under independent noise, a best case for clinical
maps (Section~\ref{sec:validation_range}), and that $\sigma$ is a width
convention: the clinical profiles decay on average more sharply than a
Gaussian, so absolute widths depend on the model, while map rankings
were strongly concordant ($\rho = 0.97$).

The width quantifies the spatial resolution of the score map in this
sense.  Its
association with the distance to the marked pacing site links this map
property to localisation performance in the present cohort: in 25 of
30 maps with $\sigma < 10$\,mm the maximum lay within 10\,mm of the
marked site, in the seven maps with $\sigma \ge 20$\,mm a median of
50\,mm away (Section~\ref{sec:clinical_localisation},
Fig.~\ref{fig:marker_check}).  In five narrow maps another location
reached the highest score, and each flag caught only one of them.  This supports a
cohort-level association between decay width and localisation offset,
not a calibrated error estimate for an individual map.
Decay width complements earlier measures of site
distinguishability~\citep{kadish1991} and of the spatial spread of the
best pacemaps~\citep{azegami2005} by giving this descriptive notion of
spatial resolution a continuous value in millimetres along the surface.

\paragraph{Distance definition.}
The distance enters the width directly: the Euclidean chord changed it
by more than 3\,mm in 16 of 55 maps, whereas exact and Heat-Method
distances agreed within 1\,mm in 48 (Section~\ref{sec:clinical_results}).
Because the synthetic fields are generated on exact geodesics and
evaluated with the Heat Method from the detected origin, the error of
the distance computation is part of the reported bias.

\paragraph{Practical use.}
A width is reported with its class and flags.  A flagged map asks for
inspection of the profile or the map, broad widths are compared only
coarsely and diffuse widths not at all, and differences between maps in
range can be judged against the precision in
Section~\ref{sec:validation_range}.  Widths should be compared under
the same convention, bin width and mapping system.

\subsection{Comparison with existing approaches}\label{sec:disc_comparison}
Table~\ref{tab:comparison} places the method next to the measures of
pace-mapping resolution used so far.  The distinguishability studies
asked whether two pacing sites a given Euclidean distance apart produce
different signals~\citep{kadish1991, azegami2005}; \citet{man1996}
derived an atrial resolution from P-wave morphology at fixed
inter-site distances; the area of the high-score region on the display
has been used as a measure of focality~\citep{lyan2024, willert2024}.
Geodesic decay analysis adds a continuous width in millimetres,
measured along the surface, with a stated operating range and a flag
for inadequate fits; like the area, it depends on the matching
algorithm of the mapping system, and unlike the distinguishability
studies it does not use an independent reference signal.  The widths of
4--6\,mm observed with the double-chamber reference are of the same
order as the ventricular distinguishability of about 5\,mm reported by
\citet{kadish1991}, and smaller than the atrial P-wave resolution of
about 17\,mm of \citet{man1996} and the 11--26\,mm over which sites with
the best ventricular pacemap score were spread~\citep{azegami2005}; but
they are a different quantity (a decay width, not a discrimination
threshold) obtained from a different signal.

The 10-mm line provides a descriptive reference.  With an amplitude of
26 score points (median in the maps with $R^2 \ge 0.7$) and a within-bin
scatter of 8 points, the fitted Gaussian predicts a score decrease at
10\,mm of 1.28 scatter units for $\sigma = 10$\,mm and 2.81 for
$\sigma = 5$\,mm.  The corresponding observed distances to the marked
pacing sites are reported in Section~\ref{sec:clinical_localisation}
(Fig.~\ref{fig:marker_check}); neither establishes a validated decision
threshold.

\begin{table*}[t]
  \centering
  \caption{Measures of spatial resolution or focality in pace mapping reported in the literature, compared with the geodesic decay width of this work.}  % values checked 26 Oct 2026 against the abstracts/full texts in method-paper/literature/
  \label{tab:comparison}
  \small\setlength{\tabcolsep}{4pt}
  \begin{tabular}{L{3.0cm}L{2.6cm}L{2.0cm}L{2.2cm}L{2.4cm}L{3.1cm}}
    \toprule
    Approach & Output & Distance & Signal / reference & Validation & Limitation \\
    \midrule
    Site distinguishability~\citep{kadish1991} & sites distinguishable at a given separation (generally within 5\,mm) & Euclidean (electrode spacing) & 12-lead ECG, ventricle & pacing at known separations & binary; Euclidean; ventricle \\
    Spread of best pacemaps~\citep{azegami2005} & greatest distance between best-score sites (11--26\,mm) & Euclidean (electroanatomical map) & 12-lead ECG, RVOT & ablation site as origin & per-patient maximum; Euclidean; ventricle \\
    P-wave resolution~\citep{man1996} & separation at which P waves change (about 17\,mm) & Euclidean (electrode spacing) & surface P wave, atrium & pacing at known separations & binary; Euclidean; surface ECG \\
    High-score area~\citep{lyan2024, willert2024} & area (cm$^2$) of the region with a similarity score $\ge 95$\,\%~\citep{willert2024} or of the best 5\,\%~\citep{lyan2024} & surface (mesh area) & intracardiac matching score & pacing from a known reference site (feasibility) & depends on the matching algorithm; no width \\
    Geodesic decay analysis (this work) & continuous width (mm), $R^2$, $F$, $\Delta\mathrm{AIC}$ & geodesic (Heat Method) & intracardiac matching score & analytic and exact geodesics; synthetic fields on clinical meshes; marked pacing sites & depends on the matching algorithm; no arrhythmia-origin or outcome validation \\
    \bottomrule
  \end{tabular}
\end{table*}

\subsection{Limitations}\label{sec:disc_limitations}
The clinical analysis used paced reference beats in 55 maps from
16 patients, analysed retrospectively, from one centre.  The localisation check was possible in
46 maps.  Prospective studies are needed to establish the relationship
between decay width, localisation of clinical arrhythmia origins and
ablation outcomes.  The reference configuration is a between-patient
factor with eight patients per group.  The data were acquired with one system (CARTO~3).  EnSite~X exports
the surface mesh and the measured point data in comparable form, and
pulse-ep imports both into the same data
structure~\citep{willert2026pulseep}, so the analysis can be applied to
such maps in the same way.  The
number of independent measurements of a map is the number of pacing
sites (7--135), not the number of mesh vertices; the synthetic
validation with independent noise at every vertex does not represent
this, and maps with few sites deserve particular caution.  In the
present maps the width was only weakly related to the number of pacing
sites (Spearman $\rho = 0.12$, 55 maps).

The model is isotropic.  Fibrosis, scar, low-voltage regions, prior
ablation lesions, local conduction abnormalities, the reference-catheter
configuration and the local anatomy can all make the decay
direction-dependent, and a high $R^2$ does not exclude anisotropy
(Section~\ref{sec:validation_shape}).  An exploratory
elliptical fit in geodesic polar coordinates (logarithmic map computed
with the Vector Heat Method~\citep{sharp2019vhm}) suggests that direction
dependence is common: in 33 maps with an interpretable fit the ratio of
the principal widths had a median of 2.7 and exceeded a within-ring
permutation null in 24 (Supplementary Section~S8).  A direction-resolved
analysis is the subject of separate work.  Because the radial Gaussian
is a width convention, it overstates the half-maximum width of sharply
peaked decays by up to 64\,\%; the per-map $\Delta\mathrm{AIC}$ and the
empirical half-maximum width are reported for this reason.

The vertex scatter of the clinical fields is larger and more
structured than the independent noise of the synthetic validation
(Section~\ref{sec:validation_shape}).  Mesh reconstruction may lengthen
geodesic distances where the reconstructed surface is rougher than the
tissue; Taubin smoothing of the meshes changed the widths by a median of
0.01\,mm or less.  A bin width scaled to each mesh, 1.5 times the mean
edge length (1.4--4.2\,mm), gave a median change in~$\sigma$ of
0.00\,mm, with five maps changing by more than 1\,mm; fixed bins of 1
and 4\,mm gave $+0.04$ and $+0.19$\,mm (Supplementary Table~S17).  The fixed
2-mm bin is kept as the default.  These sensitivities were assessed on
the meshes of this study.  The
focality index reports regions above the half maximum that are not
connected to the peak at the origin, but it does not tell a second
pacing-site response from a structured deviation of the score field.

\subsection{Future directions}\label{sec:disc_future}
A direction-resolved width can be obtained from geodesic polar
coordinates around the origin.  The estimator
applies to any peaked field on a triangulated surface; its behaviour has
been characterised only for pace-mapping score maps on atrial surfaces.  Its use
during procedures, as immediate feedback on how selective the current
map is, remains to be evaluated.


% =============================================================
%  8  CONCLUSIONS
% =============================================================
\section{Conclusions}\label{sec:conclusions}
Geodesic decay analysis quantifies the spatial resolution of a score map
as the width of its decay in millimetres, accompanied by fit-quality and focality flags: Heat-Method
geodesic distances from the vertex of maximum score, radial binning and a
Gaussian width convention, computed in under a second per map.  On
synthetic fields generated on the clinical meshes the mean bias of the
width is within $\pm 0.25$\,mm for 2--30\,mm under independent noise of
10\,\%, with a single-estimate SD of 0.2--1.8\,mm for 3--30\,mm.  Among the 46 maps with a marked pacing site, larger widths were
associated with larger distances between the score maximum and the
marked site.



% =============================================================
%  BACK MATTER
% =============================================================

\section*{Code and data availability}
The geodesic decay analysis is implemented in the open-source Python
package \texttt{pulse-ep-decay}, a module of
pulse-ep~\citep{willert2026pulseep} that can also be used on its own;
its function \texttt{analyse\_map} returns the width with its class,
$R^2$, $F$, their flags and $\Delta\mathrm{AIC}$.
It uses the open-source libraries
\texttt{potpourri3d}~\citep{sharp2021} /
\texttt{geometry-central}~\citep{sharp2019} for geodesic
computation, \texttt{pygeodesic} for the exact geodesics of the
validation and \texttt{scipy}~\citep{virtanen2020} for curve fitting.
The source code, the scripts that produce every table and figure of
this article and its supplement that is derived from data (all except
Table~\ref{tab:comparison}), the archived
result files of the validation and the per-map results of
Supplementary Table~S12 are available at
\texttt{[URL~redacted~for~review]}.  The clinical mesh and score data are
not publicly available because of data-protection restrictions; the
figures that show clinical maps can therefore not be regenerated from
the public deposit.
% <<QUERY: repository / Zenodo DOI — to be inserted at the very end (decision: Sven)>>

\section*{Funding}
This research did not receive any specific grant from funding agencies
in the public, commercial, or not-for-profit sectors.

\section*{Declaration of competing interest}
The authors declare that they have no known competing financial interests
or personal relationships that could have appeared to influence the work
reported in this paper.

\printcredits

\section*{Declaration of generative AI and AI-assisted technologies in the manuscript preparation process}
During the preparation of this work the author(s) used Claude (Anthropic) and ChatGPT (OpenAI) in order to assist with language editing and code review. After using this tool/service, the author(s) reviewed and edited the content as needed and take(s) full responsibility for the content of the published article.

\bibliographystyle{cas-model2-names}
\bibliography{references_rev1}

\end{document}
